\section{相关工作}\label{ux76f8ux5173ux5de5ux4f5c}

\section*{Related Work}

\letteredsubsection{A. 多标签图像识别}\label{a.ux591aux6807ux7b7eux56feux50cfux8bc6ux522b}

\letteredsubsection{A. Multi-Label Image Recognition}

不同于单标签，多标签图像识别允许一幅图像同时对应多个非互斥类别，因此更需要处理疾病共现、标签关联和类别不平衡等问题\hyperref[_Ref213949872]{{[}14{]}}。早期方法如Binary Relevance\hyperref[_Ref213949896]{{[}15{]}}通常将多标签任务简单拆分为多个独立二分类问题，该类方法容易忽略标签之间的相关关系，在复杂视觉场景中，往往会导致预测结果不完整或矛盾。

Unlike single-label recognition, multi-label image recognition allows one image to be assigned to multiple non-mutually exclusive categories, and therefore requires careful handling of disease co-occurrence, label relevance, and class imbalance\hyperref[_Ref213949872]{{[}14{]}}. Early methods such as Binary Relevance\hyperref[_Ref213949896]{{[}15{]}} usually decompose a multi-label task into multiple independent binary classification problems. Such methods tend to ignore label relevance and may lead to incomplete or inconsistent predictions in complex visual scenarios.

为缓解上述问题，近年来的多标签图像识别方法逐渐关注标签关联建模和样本不平衡处理。ML-GCN\hyperref[_Ref213949382]{{[}10{]}}根据疾病共现统计构建标签图，并通过图卷积网络传播标签关联信息；Asymmetric Loss (ASL)\hyperref[_Ref213949927]{{[}16{]}}对正负样本引入不同的聚焦参数，并通过负样本概率移位机制抑制易分类负样本的损失贡献，从而缓解多标签识别中的正负样本不平衡问题；Query2Label\hyperref[_Ref213950246]{{[}17{]}}进一步利用Transformer查询机制建模标签交互，并在VOC\hyperref[_Ref213950282]{{[}18{]}}等自然图像数据集上取得良好效果。这些工作表明，将标签关联显式或隐式地融入模型，有助于提升多标签预测的完整性和准确性。

To alleviate these problems, recent multi-label image recognition methods have increasingly focused on label relevance modeling and sample imbalance. ML-GCN\hyperref[_Ref213949382]{{[}10{]}} constructs a label graph from disease co-occurrence statistics and propagates label relevance information through graph convolutional networks. Asymmetric Loss (ASL)\hyperref[_Ref213949927]{{[}16{]}} introduces different focusing parameters for positive and negative samples and suppresses the loss contribution of easy negative samples through negative probability shifting, thereby alleviating positive-negative sample imbalance in multi-label recognition. Query2Label\hyperref[_Ref213950246]{{[}17{]}} further uses a Transformer query mechanism to model label interactions and achieves strong performance on natural image datasets such as VOC\hyperref[_Ref213950282]{{[}18{]}}. These studies show that explicitly or implicitly incorporating label relevance into the model can improve the completeness and accuracy of multi-label prediction.

在医疗眼底图像中，多标签识别具有更强的临床意义，因为同一只眼可能同时存在多种疾病，忽略疾病共存关系容易造成误检与漏检\hyperref[_Ref213949362]{{[}9{]}}。已有研究从不同角度探索多标签眼科疾病识别问题。例如，Cheng et al.\hyperref[_Ref213950349]{{[}19{]}}将ML-GCN迁移至眼底病灶多标签识别任务，通过构建眼底病灶语料和GloVe标签表示建模病灶标签关联；He et al.\hyperref[_Ref213950338]{{[}20{]}}提出的DCNet通过空间相关模块捕捉左右眼眼底图像特征之间的像素级关联，并融合双目特征进行患者级多标签眼科疾病识别；DKCNet\hyperref[_Ref213950343]{{[}21{]}}利用注意力模块生成判别性区域特征，并结合通道重标定与重采样策略缓解眼科数据集中的类别不平衡问题。

In medical fundus images, multi-label recognition is clinically important because multiple diseases may occur simultaneously in the same eye, and ignoring disease co-occurrence relationships may cause false positives and false negatives\hyperref[_Ref213949362]{{[}9{]}}. Existing studies have explored multi-label ocular disease recognition from different perspectives. For example, Cheng et al.\hyperref[_Ref213950349]{{[}19{]}} transferred ML-GCN to multi-label recognition of fundus lesions and modeled lesion label relevance by constructing a fundus lesion corpus and GloVe label representations. DCNet, proposed by He et al.\hyperref[_Ref213950338]{{[}20{]}}, captures pixel-level relevance between left-eye and right-eye fundus image features through a spatial correlation module and fuses binocular features for patient-level multi-label ocular disease recognition. DKCNet\hyperref[_Ref213950343]{{[}21{]}} uses an attention module to generate discriminative regional features and combines channel recalibration with resampling strategies to alleviate class imbalance in ocular datasets.

总体而言，现有方法已经能够在一定程度上建模标签关联或缓解类别不平衡，但在眼底多疾病共存场景中，疾病间关系复杂且病灶表现具有局部性，仅依赖疾病共现统计或通用注意力机制仍可能难以稳定刻画标签-特征关联性。因此，面向眼底图像构建更自适应的标签-特征关联性仍具有研究意义。

Overall, existing methods have partially addressed label relevance modeling and class imbalance. However, in fundus images with multi-disease co-occurrence, disease relationships are complex and lesion manifestations are often localized. Relying only on disease co-occurrence statistics or generic attention mechanisms may still be insufficient to characterize label-feature relevance reliably. Therefore, adaptive label-feature relevance construction for fundus images remains an important research problem.

\letteredsubsection{B. 域泛化}\label{b.ux57dfux9002ux5e94ux57dfux6cdbux5316ux4e0eux5355ux57dfux6cdbux5316}

\letteredsubsection{B. Domain Generalization}

在计算机视觉任务中，模型在训练数据上获得较好性能后，部署到新机构或新设备采集的数据上时，常会受到成像条件、光照环境和数据分布差异的影响而性能下降。这种域偏移在医疗成像中尤为突出，不同医院、设备和采集流程之间的数据差异会直接影响模型泛化能力\hyperref[_Ref213949439]{{[}22{]}}。域适应（Domain Adaptation, DA）是缓解域偏移的常见思路，其核心是在训练过程中利用源域和目标域数据进行分布对齐。例如，基于最大均值差异的Deep Domain Confusion\hyperref[_Ref213949448]{{[}23{]}}和相关性对齐方法CORAL\hyperref[_Ref213949453]{{[}24{]}}通过缩小特征分布差异提升迁移性能；DANN\hyperref[_Ref213949479]{{[}25{]}}采用对抗训练学习域不变特征；CycleGAN\hyperref[_Ref213949490]{{[}26{]}}及相关无监督域适应方法\hyperref[_Ref213949530]{{[}27{]}}则通过图像风格转换或跨域对比学习缓解域偏移。然而，DA通常要求训练阶段能够访问目标域数据，这一前提在医疗数据隐私敏感或目标机构数据不可提前获取的场景中并不总是成立。

In computer vision tasks, a model may perform well on training data but degrade when deployed on data collected by new institutions or devices because of differences in imaging conditions, illumination environments, and data distributions. This domain shift is especially prominent in medical imaging, where differences among hospitals, devices, and acquisition procedures can directly affect model generalization\hyperref[_Ref213949439]{{[}22{]}}. Domain Adaptation (DA) is a common approach for alleviating domain shift, with the core idea of aligning distributions using both source-domain and target-domain data during training. For example, maximum-mean-discrepancy-based Deep Domain Confusion\hyperref[_Ref213949448]{{[}23{]}} and the correlation alignment method CORAL\hyperref[_Ref213949453]{{[}24{]}} improve transfer performance by reducing feature distribution differences. DANN\hyperref[_Ref213949479]{{[}25{]}} uses adversarial training to learn domain-invariant features. CycleGAN\hyperref[_Ref213949490]{{[}26{]}} and related unsupervised domain adaptation methods\hyperref[_Ref213949530]{{[}27{]}} alleviate domain shift through image style transfer or cross-domain contrastive learning. However, DA usually requires access to target-domain data during training, an assumption that often fails when medical data are privacy-sensitive or data from the target institution cannot be obtained in advance.

与DA不同，域泛化（Domain Generalization, DG）旨在不访问目标域数据的条件下学习域不变特征，并泛化到未知域\hyperref[_Ref213949554]{{[}28{]}}。典型DG方法包括MixStyle\hyperref[_Ref213949562]{{[}29{]}}，通过混合实例风格统计模拟潜在域偏移；DomainBed\hyperref[_Ref213949607]{{[}30{]}}提供统一基准以评估不同DG算法；FACT\hyperref[_Ref213949788]{{[}31{]}}从频域角度分解域特定和域不变信息；RSC\hyperref[_Ref213949796]{{[}32{]}}通过抑制主导特征迫使模型关注更多与类别相关的线索；基于元学习的DG方法\hyperref[_Ref213949803]{{[}33{]}}则通过模拟源域与目标域差异来提升泛化能力。多数DG研究通常利用多个带标注源域学习跨域稳定表示，这有助于模型捕捉更丰富的域间变化，但也对标签空间统一、数据处理流程协调和跨机构协作提出较高要求。然而，在真实医疗场景中，上述条件并不总是能够满足，导致多机构数据往往难以集中或同步使用。即使通过联邦学习等方式避免集中存储，也会引入通信、隐私机制、非独立同分布数据和协同优化等额外复杂度。

Unlike DA, Domain Generalization (DG) aims to learn domain-invariant features without accessing target-domain data and to generalize to unseen domains\hyperref[_Ref213949554]{{[}28{]}}. Typical DG methods include MixStyle\hyperref[_Ref213949562]{{[}29{]}}, which simulates potential domain shift by mixing instance style statistics; DomainBed\hyperref[_Ref213949607]{{[}30{]}}, which provides a unified benchmark for evaluating DG algorithms; FACT\hyperref[_Ref213949788]{{[}31{]}}, which decomposes domain-specific and domain-invariant information from a frequency-domain perspective; RSC\hyperref[_Ref213949796]{{[}32{]}}, which suppresses dominant features to force the model to attend to more class-relevant cues; and meta-learning-based DG methods\hyperref[_Ref213949803]{{[}33{]}}, which improve generalization by simulating differences between source and target domains. Most DG studies use multiple labeled source domains to learn cross-domain stable representations. This helps models capture richer inter-domain variations, but also requires consistent label spaces, coordinated data processing, and cross-institutional collaboration. In real-world medical scenarios, however, these conditions are not always available, making it difficult to centralize or synchronously use multi-institutional data. Even when centralized storage is avoided through federated learning, additional complexity arises from communication, privacy mechanisms, non-independent and identically distributed data, and collaborative optimization.

在此背景下，单域泛化（Single Domain Generalization, SDG）作为更受限的泛化设置，仅使用一个源域学习面向未知域的域不变特征。医疗图像中的single-source DG研究\hyperref[_Ref213949863]{{[}34{]}}也说明，在无法集中使用多机构数据时，从单一源域中挖掘更稳定的疾病相关特征具有现实价值。因此，仅利用单一源域学习域不变特征，是更贴近跨机构多标签眼科疾病识别部署约束的研究方向。

Against this background, Single Domain Generalization (SDG) is a more constrained generalization setting that uses only one source domain to learn domain-invariant features for unseen domains. Studies on single-source DG in medical images\hyperref[_Ref213949863]{{[}34{]}} also indicate that mining more stable disease-related features from a single source domain has practical value when multi-institutional data cannot be centrally used. Therefore, learning domain-invariant features from a single source domain is well aligned with the deployment constraints of cross-institutional multi-label ocular disease recognition.

\letteredsubsection{C. 隐私保护}\label{c.ux9690ux79c1ux4fddux62a4}

\letteredsubsection{C. Privacy Preservation}

医疗图像和电子健康记录包含高度敏感的个人信息，模型训练、参数共享或云端推理过程均可能带来隐私泄露风险\hyperref[_Ref213950358]{{[}35{]}}。在医学人工智能应用中，隐私保护不仅关系到患者权益，也受到数据合规要求的约束。常见的隐私保护方案包括联邦学习、加密计算、安全多方计算和差分隐私等。其中，联邦学习（Federated Learning, FL）允许多个客户端在本地训练模型并上传参数或梯度，从而减少原始数据集中共享的需求\hyperref[_Ref213950376]{{[}36{]}}。例如，FedProx\hyperref[_Ref213950391]{{[}37{]}}通过添加近端项来处理客户端异构数据，缓解了协同优化不稳定问题。

Medical images and electronic health records contain highly sensitive personal information, and model training, parameter sharing, or cloud inference may all introduce privacy leakage risk\hyperref[_Ref213950358]{{[}35{]}}. In medical artificial intelligence applications, privacy preservation is tied not only to patient rights but also to data compliance requirements. Common privacy-preserving approaches include federated learning, encrypted computation, secure multi-party computation, and differential privacy. Among them, Federated Learning (FL) allows multiple clients to train models locally and upload parameters or gradients, thereby reducing the need for centralized sharing of raw data\hyperref[_Ref213950376]{{[}36{]}}. For example, FedProx\hyperref[_Ref213950391]{{[}37{]}} handles heterogeneous client data by adding a proximal term, alleviating instability in collaborative optimization.

除联邦学习外，为进一步防范隐私泄露，加密机制可以提供更强隐私保护。同态加密（Homomorphic Encryption, HE）允许在密文上执行计算\hyperref[_Ref213950427]{{[}38{]}}，并已被探索用于医疗影像分析\hyperref[_Ref213950433]{{[}39{]}}，但通常伴随较高计算开销；安全多方计算（Secure Multi-Party Computation, SMPC）\hyperref[_Ref213950438]{{[}40{]}}可在多方协作时保护各方输入隐私；可信执行环境（Trusted Execution Environment, TEE）\hyperref[_Ref213950453]{{[}41{]}}则通过硬件隔离提供更安全的计算环境。差分隐私（Differential Privacy, DP）通过向查询结果或训练梯度中注入噪声，为单个样本的隐私泄露风险提供可量化约束\hyperref[_Ref213950443]{{[}42{]}}。其中，高斯差分隐私\hyperref[_Ref213950449]{{[}43{]}}通过对裁剪后的梯度注入高斯噪声，为连续梯度更新过程提供可量化的隐私保护，因此在深度学习梯度保护中具有较强适用性。

Beyond federated learning, encryption mechanisms can provide stronger privacy preservation against privacy leakage. Homomorphic Encryption (HE) allows computation over ciphertexts\hyperref[_Ref213950427]{{[}38{]}} and has been explored for medical image analysis\hyperref[_Ref213950433]{{[}39{]}}, but it usually incurs high computational overhead. Secure Multi-Party Computation (SMPC)\hyperref[_Ref213950438]{{[}40{]}} can protect the input privacy of each party during multi-party collaboration. Trusted Execution Environment (TEE)\hyperref[_Ref213950453]{{[}41{]}} provides a more secure computing environment through hardware isolation. Differential Privacy (DP) injects noise into query results or training gradients to provide quantifiable constraints on the privacy leakage risk of individual samples\hyperref[_Ref213950443]{{[}42{]}}. Among DP mechanisms, Gaussian differential privacy\hyperref[_Ref213950449]{{[}43{]}} injects Gaussian noise into clipped gradients and provides quantifiable privacy preservation for continuous gradient updates, making it well suited to gradient protection in deep learning.

现有隐私保护机制可以降低隐私泄露风险，但在域泛化任务中仍面临额外挑战。联邦学习、同态加密和安全多方计算往往带来通信、计算或系统协同成本；差分隐私虽然部署形式相对直接，但噪声注入可能对模型精度造成影响，且这种影响在不同类别或不同数据分布上可能并不均衡\hyperref[_Ref213950458]{{[}44{]}}，这一定程度上增加了多域泛化条件下的隐私保护难度。因此，在单域泛化训练流程中探索更可控的隐私--性能平衡，对于隐私受限场景下的跨机构多标签眼科疾病识别具有实际意义。

Existing privacy-preserving mechanisms can reduce privacy leakage risk, but they still face additional challenges in domain generalization tasks. Federated learning, homomorphic encryption, and secure multi-party computation often introduce communication, computation, or system-level collaboration costs. Differential privacy is relatively straightforward to deploy, but noise injection may affect model accuracy, and this effect may vary across classes or data distributions\hyperref[_Ref213950458]{{[}44{]}}. These issues make privacy preservation more challenging under multi-domain generalization. Therefore, exploring a more controllable privacy-performance trade-off within the single domain generalization training process has practical significance for cross-institutional multi-label ocular disease recognition under privacy-restricted scenarios.
